% Part: methods
% Chapter: induction
% Section: strong-induction

\documentclass[../../../include/open-logic-section]{subfiles}

\begin{document}

\olfileid[id]{mth}{ind}{str}

\olsection{Induksi Kuat}

Dalam prinsip induksi yang dibahas di atas, kita membuktikan $P(0)$ dan juga
bahwa jika $P(k)$, maka $P(k+1)$. Pada bagian kedua, kita mengasumsikan bahwa
$P(k)$ benar dan menggunakan asumsi ini untuk membuktikan $P(k+1)$. Tentu
saja, secara ekuivalen kita dapat mengasumsikan $P(k-1)$ dan menggunakannya
untuk membuktikan $P(k)$---bagian yang penting ialah bahwa kita mampu
melakukan inferensi dari sebarang bilangan menuju penerusnya; bahwa kita
dapat membuktikan klaim yang sedang dibahas bagi sebarang bilangan dengan
asumsi bahwa klaim itu berlaku bagi pendahulunya.

Ada suatu varian prinsip induksi yang tidak hanya mengasumsikan bahwa klaim
tersebut berlaku bagi pendahulu $k-1$ dari $k$, tetapi bagi semua bilangan
yang lebih kecil daripada~$k$, lalu menggunakan asumsi ini untuk menetapkan
klaim bagi~$k$. Varian ini juga memberikan klaim $P(n)$ bagi semua~$n \in
\Nat$. Sebab, setelah menetapkan $P(0)$, dengan demikian kita telah
menetapkan bahwa $P$ berlaku bagi semua bilangan yang lebih kecil daripada
$1$. Jika kita mengetahui bahwa apabila $P(l)$ bagi semua $l<k$, maka
$P(k)$, kita khususnya mengetahui hal tersebut untuk $k=1$. Jadi, kita dapat
menyimpulkan $P(1)$. Dengan ini kita telah membuktikan $P(0)$ dan $P(1)$,
yakni $P(l)$ bagi semua $l<2$; karena kita juga memiliki kondisional bahwa
jika $P(l)$ bagi semua $l<2$, maka $P(2)$, kita dapat menyimpulkan~$P(2)$,
dan demikian seterusnya.

Sesungguhnya, jika kita dapat menetapkan kondisional umum ``untuk semua
$k$, jika $P(l)$ bagi semua $l<k$, maka $P(k)$'', kita tidak perlu lagi
menetapkan $P(0)$ secara terpisah karena hal itu sudah mengikutinya. Ingatlah
bahwa klaim umum seperti ``bagi semua $l<k$, $P(l)$'' benar jika tidak ada
$l<k$. Ini merupakan kasus kuantifikasi hampa: ``semua $A$ adalah $B$''
benar jika tidak ada $A$, dan
$\lforall[x][(!A(x) \lif !B(x))]$ benar jika tidak ada $x$ yang memenuhi
$!A(x)$. Dalam kasus ini, versi formalnya ialah
``$\lforall[l][(l < k \lif P(l))]$''---dan pernyataan itu benar jika tidak
ada $l < k$. Jika $k=0$, memang tepat itulah keadaannya: tidak ada $l<0$;
karena itu, ``bagi semua $l<0$, $P(l)$'' benar, apa pun $P$-nya. Dengan
demikian, bukti bagi ``jika $P(l)$ bagi semua $l<k$, maka $P(k)$'' secara
otomatis menetapkan~$P(0)$.

Varian ini berguna apabila penetapan klaim bagi~$k$ tidak dapat hanya
bersandar pada klaim bagi $k-1$, tetapi mungkin memerlukan asumsi bahwa
klaim itu benar bagi satu atau lebih $l<k$.

\end{document}
